\begin{theorem}[几乎所有素数阴影的刚性完备性]\label{thm:conclusion-moment-good-prime-shadow-rigidity}
固定整数 \(q\ge 2\)，令
\[
a_n:=S_q(n+2)\qquad(n\ge 0),
\]
并记其唯一首一最小整系数递推多项式为
\[
P_q(\lambda)\in \ZZ[\lambda],
\qquad
\deg P_q=d_q.
\]
设 \(H_0(q)\) 为定义 \ref{def:pom-hankel-normal-form} 中对应于该最小递推的首个可逆 \(d_q\times d_q\) Hankel 主块。则对每个满足
\[
p\nmid \det H_0(q)
\]
的素数 \(p\)，模 \(p\) 约化序列 \(\bar a_n:=a_n\bmod p\) 在 \(\FF_p\) 上仍具有同阶最小首一递推，且其最小递推多项式恰为
\[
\overline{P_q}(\lambda)=P_q(\lambda)\bmod p.
\]
因而，共有限个良素数上的最小递推阴影族
\[
\{\overline{P_q}\}_{p\nmid \det H_0(q)}
\]
唯一决定 \(P_q\) 本身；进一步，\(P_q\) 连同任意 \(d_q\) 个连续初值唯一决定整个序列 \((S_q(m))_{m\ge 2}\)。
\end{theorem}

\begin{proof}
由命题 \ref{prop:pom-moment-minrecurrence-hankel}，\(d_q\) 同时等于序列 \((a_n)\) 的最小递推阶数、Hankel 秩与最小 realization 维数。再由推论 \ref{cor:xi-hankel-minimal-recurrence-good-prime-preservation}，对每个满足
\[
p\nmid \det H_0(q)
\]
的素数，\(\bar a_n\) 在 \(\FF_p\) 上仍具有同阶最小首一递推，且其最小递推多项式正是 \(\overline{P_q}\)。

若另有首一整系数多项式 \(Q(\lambda)\neq P_q(\lambda)\) 在全部良素数上给出同一阴影，则 \(P_q-Q\) 至少有一个非零整数系数。该系数只被有限多个素数整除，故可取一良素数 \(p\nmid \det H_0(q)\) 避开这些除子，使
\[
\overline{Q}\neq \overline{P_q},
\]
与良素数上的最小性矛盾。故良素数阴影族唯一决定 \(P_q\)。

最后，由最小递推阶为 \(d_q\)，任意 \(d_q\) 个连续初值都唯一生成 \((a_n)\)，等价地唯一生成 \((S_q(m))_{m\ge 2}\)。证毕。
\end{proof}

\begin{corollary}[有限良素数即可作整数级重建]\label{cor:conclusion-moment-finite-good-primes-crt-reconstruction}
在定理 \ref{thm:conclusion-moment-good-prime-shadow-rigidity} 的记号下，若事先给定一个对 \(P_q\) 全部系数与所选 \(d_q\) 个连续初值的统一高度上界 \(B_q\)，并取若干良素数
\[
p_1,\dots,p_r\nmid \det H_0(q),
\qquad
\prod_{i=1}^{r}p_i>2B_q,
\]
则仅凭这些模 \(p_i\) 的阴影递推与阴影初值，即可经中国剩余定理唯一恢复 \(P_q\) 与相应初值，从而恢复全部 \((S_q(m))_{m\ge 2}\)。
\end{corollary}

\begin{proof}
由定理 \ref{thm:conclusion-moment-good-prime-shadow-rigidity}，每个良素数 \(p_i\) 上的阴影都给出真递推多项式 \(P_q\) 与所选初值块的模 \(p_i\) 余数。于是中国剩余定理给出模
\[
P:=\prod_{i=1}^{r}p_i
\]
的唯一剩余类。由于 \(P>2B_q\)，每个绝对值不超过 \(B_q\) 的整数只能对应唯一的整代表，故 \(P_q\) 的系数与全部初值都被唯一恢复。最后由递推唯一性恢复整个序列。证毕。
\end{proof}

\begin{proposition}[坏素数集即 Hankel 刚性支撑]\label{prop:conclusion-moment-hankel-bad-prime-stiffness-support}
在定理 \ref{thm:conclusion-moment-good-prime-shadow-rigidity} 的设定下，可能导致递推阶下降、最小多项式变形或最小 realization 算术失真的素数集必包含于
\[
\operatorname{Supp}\bigl(\det H_0(q)\bigr).
\]
更精确地，对每个素数 \(p\) 定义
\[
\operatorname{Stiff}_p(q):=v_p\!\bigl(\det H_0(q)\bigr),
\]
并记
\[
H_1(q):=\bigl(a_{i+j+1}\bigr)_{0\le i,j\le d_q-1}
\]
为与 \(H_0(q)\) 相邻的偏移 Hankel 块，则 \(\operatorname{Stiff}_p(q)\) 同时控制：
\begin{enumerate}
  \item 模 \(p\) 阴影是否可能失真；
  \item 最小有理 realization
  \[
  A=H_0(q)^{-1}H_1(q)
  \]
  的 \(p\)-进恢复门槛；
  \item 整数递推与 \(p\)-进最小模型之间的算术刚性。
\end{enumerate}
\end{proposition}

\begin{proof}
由推论 \ref{cor:xi-hankel-minimal-recurrence-good-prime-preservation}，当 \(p\nmid \det H_0(q)\) 时，模 \(p\) 阴影保持最小递推阶与最小首一多项式；故任何模 \(p\) 失真都只能发生在
\[
p\mid \det H_0(q)
\]
时。

另一方面，最小 Hankel realization 的 companion-normal form 由
\[
A=H_0(q)^{-1}H_1(q)
=\frac{\operatorname{adj}(H_0(q))\,H_1(q)}{\det H_0(q)}
\]
给出，因此其全部有理分母都整除 \(\det H_0(q)\)。故 \(p\)-进恢复精度与算术稳定性只可能在
\[
v_p\!\bigl(\det H_0(q)\bigr)>0
\]
时出现损失，而其损失强度恰由该赋值 \(\operatorname{Stiff}_p(q)\) 控制。证毕。
\end{proof}

\begin{theorem}[连续压力的整数阶代数锚点]\label{thm:conclusion-projective-pressure-integer-algebraic-anchor}
对每个整数 \(q\ge 2\)，射影压力算子 \(\mathcal L_q\) 在齐次次数 \(q\) 子空间 \(\mathcal H_q\) 上满足严格代数退化
\[
\mathcal L_q\vert_{\mathcal H_q}=\widetilde A_q,
\qquad
\rho(\mathcal L_q\vert_{\mathcal H_q})=\rho(\widetilde A_q)=r_q.
\]
若连续压力函数按注 \ref{rem:pom-projective-pressure-route} 定义为
\[
\Lambda(t):=\log\lambda(t+1)-\log |A_{\mathrm{out}}|,
\]
其中 \(\lambda(t+1)\) 是 \(\mathcal L_{t+1}\) 的 Perron 主特征值，则在所有整数锚点 \(t=q-1\) 处都有精确恒等式
\[
\Lambda(q-1)=\log r_q-\log |A_{\mathrm{out}}|.
\]
对于 Fold 主核，\(|A_{\mathrm{out}}|=2\)，故
\[
\Lambda(q-1)=\log r_q-\log 2.
\]
\end{theorem}

\begin{proof}
定理 \ref{thm:pom-projective-operator-degenerates-to-moment-kernel} 已证明：对每个整数 \(q\ge 2\)，\(\mathcal L_q\) 在 \(\mathcal H_q\) 上的限制在任一单项式基下恰由 \(\widetilde A_q\) 表示，且
\[
\rho(\mathcal L_q\vert_{\mathcal H_q})=\rho(\widetilde A_q)=r_q.
\]
再由注 \ref{rem:pom-projective-pressure-route} 对连续压力的定义，整数点 \(t=q-1\) 正是该有限维退化的精确插值节点，因此
\[
\Lambda(q-1)=\log\lambda(q)-\log |A_{\mathrm{out}}|
=\log r_q-\log |A_{\mathrm{out}}|.
\]
对 Fold 主核代入 \(|A_{\mathrm{out}}|=2\) 即得最后一式。证毕。
\end{proof}

\begin{theorem}[大偏差骨架由离散湮灭多项式族完全承载]\label{thm:conclusion-large-deviation-skeleton-carried-by-annihilator-family}
令
\[
\mathscr P:=\{P_q(\lambda)\}_{q\ge 2}
\]
为全部碰撞矩最小递推多项式族。则 \(\mathscr P\) 连同相应最小初值块，唯一决定全部整数锚点压力值
\[
\{\Lambda(q-1)\}_{q\ge 2}.
\]
换言之，连续大偏差接口在整数方向上的全部可核验内容，已被离散湮灭多项式族 \(\mathscr P\) 完整编码。
\end{theorem}

\begin{proof}
对每个固定 \(q\ge 2\)，由定理 \ref{thm:conclusion-moment-good-prime-shadow-rigidity}，\(P_q\) 与其最小初值块唯一决定整个整数序列 \((S_q(m))_{m\ge 2}\)。再由 Perron--Frobenius 理论，任一有限维非负 realization 的主增长率唯一决定于该序列的指数尺度，故 \((S_q(m))\) 唯一决定 \(r_q\)。最后由定理 \ref{thm:conclusion-projective-pressure-integer-algebraic-anchor}，
\[
\Lambda(q-1)=\log r_q-\log 2
\]
即被唯一确定。故整数锚点压力值族完全由 \(\mathscr P\) 承载。证毕。
\end{proof}
